\subsection{维特分解。}

除了本小节开头以来已经生效的约定之外，我们再补充如下约定：

条件 (T). --- 对所有 \(x \in E\)，存在 \(\alpha \in A\) 使得 \(\Phi(x, x) = \alpha + \varepsilon\bar{\alpha}\)。

当 \(\Phi\) 是交错型，或当 \(\varepsilon = 1\) 且 A 是特征 \(\neq 2\) 的体时，取 \(\alpha = \frac{1}{2}\Phi(x, x)\)，该条件总是满足的（参见习题 1 与 14）。

引理 1. --- 设 \(\Phi\) 是 E 上满足 (T) 的 \(\varepsilon\)-埃尔米特型（相应地，与二次型 Q 相伴的双线性型），F 是 E 的不约化为 0 的全迷向（相应地，全奇异）子空间。对每个不与 F 正交的 \(x \in E\) 和所有 \(\alpha \in A\)，存在 \(y \in F\) 使得

\[
\Phi (x + y, x + y) = \alpha + \varepsilon \bar {\alpha} \quad (\text { resp. } Q (x + y) = \alpha).
\]

事实上，设 \(\Phi(x, x) = \beta + \varepsilon\bar{\beta}\)（相应地，\(Q(x) = \beta\)）。那么对 \(y \in F\)，有 \(\Phi(x + y, x + y) = (\beta + \Phi(x, y)) + \varepsilon(\overline{\beta + \Phi(x, y)})\)（由于 \(\Phi(y, y) = 0\)）（相应地，有 \(Q(x + y) = \beta + \Phi(x, y)\)，由于 \(Q(y) = 0\)）。由于 \(x\) 不与 \(F\) 正交，仿射线性函数 \(y \to \Phi(x, y) + \beta\) （\(F\) 上的）不是常数；因此它取值 \(\alpha\)（对某个元素 \(y\)，该元素属于 \(F\)），于是即回答了问题。

E 的维特分解是指将 E 分解为三个子空间 F、F'、G 的直和，使得 F 与 F' 全迷向（相应地，全奇异）且 G 非迷向并与 F + F' 正交；若 E 是有限维的，则 Φ 关于适应于 E 的维特分解的 E 的基的矩阵具有形式

\[
\left( \begin{array}{c c c} 0 & U & 0 \\ \varepsilon \overline {{U}} & 0 & 0 \\ 0 & 0 & V \end{array} \right)\tag{1}
\]

我们称 \(\Phi\) 为中性型，如果它非退化，且 E 是有限维的并且是两个全迷向（相应地，全奇异）子空间的直和。两个中性型的直和是中性型。

命题 2. --- 设 Φ 是满足 (T) 的非退化 ε-埃尔米特型（相应地，与非退化二次型 Q 相伴的双线性型），F 是有限维 r 的全迷向（相应地，全奇异）子空间。

\begin{enumerate}
\def\labelenumi{\alph{enumi})}
\item
  若 \(F'\) 是维数为 r 且满足 \(F' \cap F^{0} = \{0\}\) 的全迷向子空间，则 \(F + F'\) 非迷向，且对 F 的任意基 \((f_{i})\)，存在基 \((f_{i}')\) （\(F'\) 的），使得 \(\Phi(f_{i}, f_{j}') = \delta_{ij}\)（克罗内克指标）对 \(i, j = 1, \ldots, r\) 成立。
\item
  若 G 是维数 \(\leqslant r\) 且满足 \(G \cap F^{0} = \{0\}\) 的全迷向（相应地，全奇异）子空间，则存在全迷向（相应地，全奇异）子空间 \(F' \supset G\)，维数为 r 且满足 \(F' \cap F^{0} = \{0\}\)。
\end{enumerate}

设 \(\Psi\) 是 \(\Phi\) 在 \(F \times F'\) 上的限制；对 \(x' \in F'\)，关系 \(\langle \Phi(x, x') = 0\) 对所有 \(x \in F\rangle\) 蕴含 \(x = 0\)，因为 \(F' \cap F^0 = \{0\}\)。于是断言 \(a)\) 由 §1 第 6 小节命题 6 的推论得出，只是 \(F + F'\) 非迷向这一点除外。而子空间 \(H = (F + F') \cap (F + F')^0\) 等于 \((F + F') \cap F^0 \cap F'^0\)。由于 \(F \subset F^0\)，有 \((F + F') \cap F^0 = F + (F' \cap F^0) = F\)，从而 \(H = F \cap F'^0\)；故 \(H = \{0\}\)（既然我们已看到 \(\Psi\) 非退化）。这就证明了 \(F + F'\) 非迷向。

为证明 \(b\) )，我们对 \(s = \dim G\) 作递降归纳。因此只需证明：若 \(s < r\)，则存在全迷向（相应地，全奇异）子空间 \(G'\)，它包含 \(G\)、维数为 \(s + 1\) 且满足 \(G' \cap F^0 = \{0\}\)。由于 \(\dim G < \dim F\)，\(\Phi\) 在 \(F \times G\) 上的限制是退化的，又因 \(G \cap F^0\) 为零，故 \(F \cap G^0\) 非零。若此时有 \(G + F^0 \supset G^0\)，则取正交子空间并注意 \(F = F^{00}\) 与 \(G = G^{00}\)（ \(\S 1\) ，第 6 小节，命题 4 的推论 1），便可推出 \(G^0 \cap F \subset G\)，从而

\[
\mathrm{G} ^ {0} \cap \mathrm{F} \subset \mathrm{G} \cap \mathrm{F} \subset \mathrm{G} \cap \mathrm{F} ^ {0} = \{0 \},
\]

这是不可能的。于是存在元素 \(x\) （\(G^0\) 的）使得 \(x \notin G + F^0\)；由于 \(F \subset F^0\)，可以给 \(x\) 加上 \(G^0 \cap F\) 的一个向量而不改变这些性质；由于 \(G^0 \cap F\) 全迷向（相应地，全奇异）且 \(\neq \{0\}\)，引理 1 表明可以选择 \(x\) 是迷向的（相应地，奇异的）。

于是子空间 \(G' = G + Ax\) 的维数为 \(s + 1\)，且是全迷向的（相应地，全奇异的）；此外我们有 \(G' \cap F^{0} = \{0\}\)，因为若 \(y = z + ax\)（\(z \in G, a \in A\)）属于 \(F^{0}\)，则有 a = 0，否则将有 \(x \in F^{0} + G\)，与 x 的选取相矛盾，从而 \(y = z \in G \cap F^{0} = \{0\}\)，即 y = 0。因此子空间 \(G'\) 回答了问题。

推论 1. --- 若 F 是维数为 r 的全迷向（相应地，全奇异）子空间，则存在维数为 r 的全迷向（相应地，全奇异）子空间 F'，使得 F ∩ F' = \{0\} 且 F + F' 非迷向。

在命题 2 b) 中取 G = \{0\} 即可。

推论 2. --- A 上相同维数的空间上的两个中性 ε-埃尔米特型是等价的。

注. --- 在推论 1 的条件下，E 是 F + F' 与 F + F' 的正交子空间的直和。这样就得到 E 的一个维特分解。由命题 2 a)，存在 F 与 F' 的基，使得在 Φ 的矩阵 (1) 中，块 U 是单位矩阵 1。

命题 3. --- 设 \(\Phi\) 是满足 (T) 的非退化 \(\varepsilon\)-埃尔米特型（相应地，与非退化二次型 Q 相伴的双线性型）。设 \(F_{1}\) 与 \(F_{2}\) 是 E 的两个极大全迷向（相应地，全奇异）子空间，其中一个是有限维的。令 \(F = F_{1} \cap F_{2}\)。设 \(S_{i}\)（\(i = 1, 2\)）是 F 在 \(F_{i}\) 中的补；令 \(S = S_{1} + S_{2}\)。则存在 E 的两个子空间 G 和 H，使得

\begin{enumerate}
\def\labelenumi{\alph{enumi})}
\item
  子空间 G + F、S 与 H 非迷向且两两正交；
\item
  E 是 F、S、G 与 H 的直和；
\item
  H 中没有非零的迷向（相应地，奇异）向量；
\item
  G 是全迷向的（相应地，全奇异的）。
\end{enumerate}

此外，\(F_{1}\) 与 \(F_{2}\) 都是有限维的，并且有 \(\dim \mathrm{F}_1 = \dim \mathrm{F}_2, \dim \mathrm{G} = \dim \mathrm{F}, \dim \mathrm{S}_1 = \dim \mathrm{S}_2, \operatorname{codim} \mathrm{H} = 2 \dim \mathrm{F}_1.\)

首先注意，若 N 是极大全迷向（相应地，全奇异）子空间，则每个与 N 正交的迷向（相应地，奇异）向量 x 都属于 N，因为否则 \(N + Ax\) 将与 N 的极大性矛盾。于是，若对 i = 1, 2，\(x_{i}\) 是 \(F_{i}^{0}\) 的迷向（相应地，奇异）向量，则有 \(x_{i} \in F_{i}\)。另一方面，若 y 是 \(S_{1}\) 中与 \(S_{2}\) 正交的元素，则由于 \(F_{1}\) 全迷向，y 与 \(F_{1}\) 正交，从而与 F 正交，进而与 \(F_{2} = S_{2} + F\) 正交。由于 y 是迷向（相应地，奇异）的且与 \(F_{2}\) 正交，故有

\[
y \in S _ {1} \cap F _ {2} = S _ {1} \cap F _ {1} \cap F _ {2} = S _ {1} \cap F = \{0 \}.
\]

于是有 \(S_1 \cap S_2^0 = \{0\}\)，同样 \(S_2 \cap S_1^0 = \{0\}\)。由于两个子空间 \(F_1, F_2\) 之一（例如 \(F_1\)）是有限维的，故 \(S_1\) 有限维，从而 \(S_1^0\) 有有限余维数（§ 1，第 6 小节，命题 4 的推论 1），于是 \(S_2\) 也是有限维的，因为 \(S_2 \cap S_1^0 = \{0\}\)；此外这表明 dim \(S_2 \leqslant \text{codim } S_1^0 = \text{dim } S_1\)；同样 dim \(S_1 \leqslant \text{dim } S_2\)，从而 dim \(S_1 = \text{dim } S_2\)。于是命题 2 a) 表明 \(S = S_1 + S_2\) 非迷向。

在此基础上，S 的正交子空间 N 非迷向（第 1 小节，命题 1 的推论）且包含 F；因此命题 2 的推论 1 表明，存在 N 的全迷向（相应地，全奇异）子空间 G，使得 dim G = dim F、G ∩ F = \{0\} 且 G + F 非迷向。于是 G 满足 d)。再取 H 为 G + F 在 N 中的正交子空间，便满足 a) 与 b)。至于 c)，注意由于 H 与 F₁ = S₁ + F 正交，根据证明开头所作的说明以及 H ∩ F₁ = \{0\} 这一事实，H 中没有非零的迷向（相应地，奇异）向量。最后，关于维数的某些断言已在证明过程中得出；其余的由它们平凡地推出。

推论 1. --- 在命题 3 的假设下，两个有限维的极大全迷向（相应地，全奇异）子空间有相同的维数。对每个有限维的极大全迷向（相应地，全奇异）子空间 F，存在另一个这样的子空间 \(\mathbf{F}'\)，使得 \(\mathbf{F} \cap \mathbf{F}' = \{0\}\)，且在这些条件下 \(\mathbf{F} + \mathbf{F}'\) 非迷向。

若 \(F \cap F' = \{0\}\)，则按命题 3 的记号将有 \(G = \{0\}\)，且 \(F + F'\) 非迷向。其余断言由命题 3 与命题 2 的推论 1 平凡地得出。

推论 2. --- 设 Q 是代数闭体 A 上有限维 n 的向量空间 E 上的非退化二次型；则存在 E 的基 \((e_i)_{1 \leqslant i \leqslant n}\)，使得

\begin{enumerate}
\def\labelenumi{(\arabic{enumi})}
\setcounter{enumi}{1}
\tightlist
\item
\end{enumerate}

\[
\mathrm{Q} (\sum_ {i = 1} ^ {n} x _ {i} e _ {i}) = \sum_ {i = 1} ^ {\nu} x _ {i} x _ {i + \nu} \quad \text {   if   } n = 2 \nu ,\tag{3}
\]

\[
\mathrm{Q} \left(\sum_ {i = 1} ^ {n} x _ {i} e _ {i}\right) = \sum_ {i = 1} ^ {\nu} x _ {i} x _ {i + \nu} + x _ {2 \nu + 1} ^ {2} \quad \text {   if   } n = 2 \nu + 1.
\]

事实上，设 \(F_{1}\) 与 \(F_{2}\) 是满足 \(F_{1} \cap F_{2} = \{0\}\)（推论 1）的两个极大全奇异子空间，q 是它们的维数。于是按命题 3 的记号有 \(G = \{0\}\)。取基 \((e_{i})_{1 \leqslant i \leqslant q}\) （\(F_{1}\) 的）与基 \((e_{i})_{q+1 \leqslant i \leqslant 2q}\) （\(F_{2}\) 的），后者满足 \(\Phi(e_{i}, e_{j+q}) = \delta_{ij}\)（对 \(i, j = 1, \ldots, q\)）（命题 2 a)），我们看出只需证明 dim \(H \leqslant 1\)。而若 \(x \in H\)、\(y \in H\) 且 \(x \neq 0\)，则方程 \(Q(y - ax) = Q(y) - a\Phi(x, y) + a^{2}Q(x) = 0\) 至少有一个解 \(a_{0}\)（由于 \(Q(x) \neq 0\)），又由于 H 的每个奇异向量都是零，故有 \(y = a_{0}x\)。

定义 3. --- 假设 E 是有限维的，且 \(\Phi\) 是满足 (T) 的非退化 \(\varepsilon\)-埃尔米特型（相应地，与非退化二次型 Q 相伴的双线性型）。\(\Phi\)（相应地，Q）的指数是 E 的极大全迷向（相应地，全奇异）子空间的公共维数。

若 n 是 E 的维数，v 是 \(\Phi\)（相应地，Q）的指数，则命题 3 表明有

\[
n \geqslant 2 v.\tag{4}
\]

此外，由于每个全迷向（相应地，全奇异）子空间都包含在某个极大全迷向（相应地，全奇异）子空间之中，故极大全迷向（相应地，全奇异）子空间恰好是维数为 v 的那些。\(\Phi\)（相应地，Q）的指数为 0 这一断言意味着 E 的每个迷向（相应地，奇异）向量都是零。在偶数维 \(n\) 的空间中，中性型恰好是指数为 \(\frac{1}{2} n\) 的那些；在奇数维空间中没有中性型。命题 3 表明，每个型都是一个中性型与一个指数为 0 的型的直和。

命题 4. --- 设 Q 是 E 上的非退化二次型，且存在 E 的向量 \(x \neq 0\) 使得 Q(x) = 0。则对 A 的每个元素 a，存在 \(y \in E\) 使得 Q(y) = a。

事实上，由命题 2 的推论 1，存在子空间 \(\mathbf{G} = \mathbf{F} + \mathbf{F}'\)（\(\mathbf{F}, \mathbf{F}'\)：维数为 1 的全奇异子空间；\(\mathbf{E}\) 的、维数为 2 的），使得 \(\mathbf{Q}\) 在 \(\mathbf{G}\) 上的限制是中性的。若 \(\{e, e'\}\)（\(e \in \mathbf{F}, e' \in \mathbf{F}'\)）是 \(\mathbf{G}\) 的基，则有

\[
\mathrm{Q} (x e + x ^ {\prime} e ^ {\prime}) = b x x ^ {\prime} \quad (x \in \mathrm{A}, x ^ {\prime} \in \mathrm{A}, b \in \mathrm{A}, b \neq 0).
\]

因此取 \(ae + b^{-1}e'\) 为 y 即可。
% semantic-fragments: legacy-input-seam
\relax{}
